% Introductory supervised-learning lab (TASKS.md D3: short guided PyTorch lab); faculty draft.
% Shared layout adapted from the instructor's ME 537 homework sources.
\documentclass[12pt,letterpaper]{article}
\usepackage{mathptmx}
\usepackage[T1]{fontenc}
\usepackage{xcolor,amsmath,amssymb,xurl,listings,enumitem,tikz}
\textwidth=6.5in
\textheight=9in
\oddsidemargin=0in
\evensidemargin=0in
\topmargin=0in
\headheight=0in
\headsep=0in
\parskip=0.08in
\parindent=0in
\setlist[enumerate]{itemsep=0.04in,topsep=0.04in}
\lstset{basicstyle=\ttfamily\footnotesize,breaklines=true,columns=fullflexible,
keepspaces=true,showstringspaces=false,aboveskip=0.04in,belowskip=0.04in}
\newcommand{\R}[2]{{}^{#1}\!R_{#2}}
\newcommand{\T}[2]{{}^{#1}\!T_{#2}}
\newcommand{\courseheader}[2]{%
\setlength{\fboxrule}{2pt}\setlength{\fboxsep}{2mm}%
\fcolorbox{black}{white}{\parbox{\dimexpr\linewidth-2\fboxsep-2\fboxrule\relax}{\centering\textbf{Introduction to Robotics Homework \##1 -- Draft}}}\par
\textbf{#2}\par}
\newcommand{\conventions}{Use right-handed rotations, column vectors, radians,
and meters. $\R{A}{B}$ maps coordinates from frame $B$ into frame $A$;
$\T{A}{B}$ also includes the origin of $B$ expressed in $A$.}
\newcommand{\submission}{\textbf{Submit:} labeled hand derivations and sketches,
your changed Python source files, the requested numerical/visual evidence,
and an AI-use statement (including ``none'' if applicable). Show your reasoning
before using code to check hand calculations. Public checks are intentionally
limited. The instructor supplies the deadline and submission destination.}

\begin{document}
\courseheader{10}{Learning to classify when rules fail}
\textbf{Schedule:} week 10, after the two supervised-learning lectures.
About three to four hours, including a training run of a minute or two on
your CPU.

\textbf{Task.} Your HW08 rule classifier (color threshold, contours, corner
count) is perfect on the project's images. Real cameras are less kind: light
gets brighter or dimmer and changes color, objects blur, noise grows, and an
arm link or cable can cover part of an object. You will train a small
convolutional network on $48\times48$ crops of single objects from the
project's fixed camera, under both \emph{clean} (project) and \emph{hard}
conditions, and compare it with your rules on scenes neither has seen. Each
crop is one of 9 classes: a color (red, blue, yellow) and a shape (circle,
square, triangle). The data are generated from fixed seeds, so everyone has
the same crops. Your code goes in \path{homework/hw10/learning.py}.

\begin{enumerate}
\item \textbf{Look at the data first.}
\begin{lstlisting}
python -m homework.hw10.hw10_data --save tmp/hw10_data.png
\end{lstlisting}
\begin{enumerate}
\item Report the number of crops and scenes in each split, the image
array's shape and type, and the shape, type, and value range of a batch
after \texttt{to\_tensor}. What does each of the four tensor dimensions mean?
\item A \emph{scene} is one object in 6 consecutive frames under one set of
conditions, so its crops are nearly identical. The splits are made by scene,
never by crop. If crops of one scene were in both the training and test
sets, would the test accuracy be too high or too low? Why?
\item Name three hard conditions in your figure and the step of your HW08
classifier (threshold, contour, or corner count) that each one would break.
\end{enumerate}

\item \textbf{The training step.}
\begin{enumerate}
\item The network gives 9 scores (logits) per crop; the cross-entropy loss is
$-\ln p_y$, where $p_k=e^{s_k}/\sum_j e^{s_j}$ and $y$ is the true class. By
hand, for three classes with scores $s=(2,\ 1,\ -1)$, find $p$ and the loss
when the true class is 0 and when it is 2. Why does a confident wrong answer
cost so much more than an unsure one?
\item Complete \texttt{train\_epoch}: the four lines that compute the loss,
clear old gradients, backpropagate, and step the optimizer, then record the
loss. What does \texttt{loss.backward()} compute, and what would go wrong
without \texttt{optimizer.zero\_grad()}?
\item Run \texttt{python -m homework.hw10.hw10\_train} (5 epochs). Report each
epoch's loss and validation accuracy (clean and hard), the total time, and
your plot (\path{tmp/hw10/cnn.png}). With 64 crops per batch, how many
optimizer steps is one epoch? Why judge the network on validation scenes
rather than on the training loss?
\end{enumerate}

\item \textbf{Evaluate against your rules.}
\begin{enumerate}
\item Complete \texttt{predict}. Why must it call \texttt{model.eval()} and
use \texttt{torch.no\_grad()}?
\item Complete \texttt{confusion\_matrix}, then run
\begin{lstlisting}
python -m homework.hw10.hw10_evaluate --save tmp/hw10_eval.png
\end{lstlisting}
It runs your network and your HW08 \texttt{find\_objects} on the same 600
test scenes. Report both classifiers' accuracy on clean and on hard crops,
and include the two confusion matrices.
\item Your rules almost never get the \emph{color} wrong on hard crops, yet
often get the shape wrong or find nothing. Explain both using your answers
to 1(c). What do the network's mistakes (bottom of the figure) have in
common?
\item For the project as it is (S1--S3, clean images), which classifier
would you use, and why? Name one change to the project that would make the
network worth its cost. What is missing from these synthetic crops before
you would trust the network on a real camera?
\end{enumerate}
\end{enumerate}

\textbf{If training is too slow} on your computer, evaluate the supplied
network, trained the same way, and say so in your answers:
\begin{lstlisting}
python -m homework.hw10.hw10_evaluate --checkpoint homework/hw10/checkpoints/cnn.pt
\end{lstlisting}

\textbf{Project connection:} a project team may use either classifier; the
evaluation's images are clean. HW11 trains a policy network with the same
PyTorch pieces.

\submission
\end{document}
